\documentclass[11pt]{letter}
\usepackage[margin=1in]{geometry}
\usepackage{times}
\usepackage{xcolor}

\signature{Jeff Gore (Corresponding Author), Jinyeop Song, and Jiliang Hu \\ Department of Physics \\ Massachusetts Institute of Technology \\ Cambridge, MA, USA \\ jeffgore@gmail.com}

\address{Department of Physics \\ Massachusetts Institute of Technology \\ Cambridge, MA 02139, USA}

\begin{document}

\begin{letter}{}

\opening{Dear Editor,}

We are pleased to submit a revised version of our Article, ``Interspecies Interactions Drive Community-Level Selection in Microbial Coalescence,'' for continued consideration in Nature Ecology \& Evolution. We are grateful to the editor and reviewers for their careful evaluation of the manuscript. Their comments helped us sharpen the conceptual framing, add important controls, and state the scope of our conclusions more precisely.

The revision addresses the reviewers' major concerns in three main ways. First, we now define community-level selection in operational terms, as origin-correlated persistence during coalescence, and quantify it using parental-community similarity, pairwise selection-correlation analyses, and fate-concordance analyses. This framing separates the outcome-level pattern from the causal ecological or biochemical mechanisms, which may vary across systems. Second, we clarify the relationship between model and experiment by distinguishing the interaction-strength parameter in the phenomenological generalized Lotka--Volterra model from experimentally measured invasion resistance, assay-based proxies for net interaction intensity, and pH-associated environmental feedbacks across nutrient conditions. Third, we expanded the analyses of alternative explanations, including classification-boundary sensitivity, additive null models, richness and abundance-imbalance controls, pH and dominant-species controls, interaction-matrix extensions, and taxonomic-distinctness analyses for communities derived from natural samples.

We also revised the figures, Extended Data, Supplementary Information, and figure captions to make the assumptions and interpretations clearer. In particular, Fig.~2A now clarifies the interaction-matrix structure used in the model workflow, while Fig.~2C and the Supplementary Information show how assembly alters within-community interaction structure. Together, these revisions clarify the central conclusion and make the mechanistic interpretation and scope more precise: assembly-generated interaction structure, effective ecological coupling, and environmental feedbacks can shape whether microbial coalescence proceeds largely as independent species sorting or whether parental-community origin becomes predictive of persistence.

Changes in the revised manuscript and Supplementary Information are marked in red for ease of review. We have also submitted a detailed point-by-point response to all reviewer comments as a separate response-to-reviewers file.

The revised Article contains six main figures and no tables. The compiled revised manuscript is 28 pages, including Methods, references, and figure legends. The revised Supplementary Information is 57 pages and includes Extended Data Figs.~1--8 and Supplementary Figs.~1--46.

We confirm that this manuscript has not been published elsewhere and is not under consideration by another journal. No related manuscripts are under consideration or in press elsewhere. All authors have approved the revised manuscript and agree with its resubmission to Nature Ecology \& Evolution. We have no competing interests to disclose. Correspondence should be addressed to Jeff Gore.

Thank you for considering our revised manuscript.

\closing{Sincerely,}

\end{letter}

\end{document}
